% GENERATED FILE - do not hand-edit.
% Source of truth: data/calculus-reviewed/t-7.json in the zenith-website repo.
% Regenerate with: python3 scripts/build_calculus_tex.py
\documentclass[11pt]{article}
\usepackage[utf8]{inputenc}
\usepackage[T1]{fontenc}
\usepackage{amsmath,amssymb}
\usepackage{graphicx}
\usepackage[margin=1in]{geometry}
\graphicspath{{../..//}}
\setlength{\parindent}{0pt}
\setlength{\parskip}{0.6em}
\begin{document}
\section*{Chapter 7: Application of Integration to Geometry --- Test T7}

\bigskip\noindent\textbf{T1.}\ Find the area of the region with the following boundaries: The curve $y=x^2$, $y=0$, $x=-2$, and $x=3$.
\begin{list}{}{\setlength{\leftmargin}{3em}\setlength{\labelwidth}{2.5em}\setlength{\itemsep}{0pt}\setlength{\topsep}{2pt}}
\item[(A)] $\dfrac{11}{3}$
\item[(B)] $\dfrac{35}{3}$
\item[(C)] $5$
\item[(D)] $7$
\end{list}

\bigskip\noindent\textbf{T2.}\ Find the area of the region with the following boundaries: The parabola $y=x^2-5$ and the line $y=4$.
\begin{list}{}{\setlength{\leftmargin}{3em}\setlength{\labelwidth}{2.5em}\setlength{\itemsep}{0pt}\setlength{\topsep}{2pt}}
\item[(A)] $9$
\item[(B)] $18$
\item[(C)] $36$
\item[(D)] $72$
\end{list}

\bigskip\noindent\textbf{T3.}\ Find the area of the region with the following boundaries: The curve $x=y^2-4$ and the $y$-axis.
\begin{list}{}{\setlength{\leftmargin}{3em}\setlength{\labelwidth}{2.5em}\setlength{\itemsep}{0pt}\setlength{\topsep}{2pt}}
\item[(A)] $\dfrac{32}{3}$
\item[(B)] $\dfrac{16}{3}$
\item[(C)] $8$
\item[(D)] $4$
\end{list}

\bigskip\noindent\textbf{T4.}\ Find the area of the region with the following boundaries: The parabola $y^2=x$ and the line $x+y=6$.
\begin{list}{}{\setlength{\leftmargin}{3em}\setlength{\labelwidth}{2.5em}\setlength{\itemsep}{0pt}\setlength{\topsep}{2pt}}
\item[(A)] $\dfrac{45}{4}$
\item[(B)] $\dfrac{125}{6}$
\item[(C)] $\dfrac{25}{2}$
\item[(D)] $\dfrac{45}{2}$
\end{list}

\bigskip\noindent\textbf{T5.}\ Find the area of the region with the following boundaries: The curve $y=\dfrac{6}{x^2+9}$, the $x$-axis, and the vertical lines $x=-3\sqrt{3}$ and $x=3\sqrt{3}$.
\begin{list}{}{\setlength{\leftmargin}{3em}\setlength{\labelwidth}{2.5em}\setlength{\itemsep}{0pt}\setlength{\topsep}{2pt}}
\item[(A)] $\dfrac{\pi}{3}$
\item[(B)] $\dfrac{2\pi}{3}$
\item[(C)] $\dfrac{4\pi}{3}$
\item[(D)] $\dfrac{8\pi}{3}$
\end{list}

\bigskip\noindent\textbf{T6.}\ Find the area of the region with the following boundaries: The parabolas $x=y^2-6y$ and $x=4y-y^2$.
\begin{list}{}{\setlength{\leftmargin}{3em}\setlength{\labelwidth}{2.5em}\setlength{\itemsep}{0pt}\setlength{\topsep}{2pt}}
\item[(A)] $\dfrac{25}{3}$
\item[(B)] $\dfrac{125}{6}$
\item[(C)] $\dfrac{125}{3}$
\item[(D)] $\dfrac{250}{3}$
\end{list}

\bigskip\noindent\textbf{T7.}\ Find the area of the region with the following boundaries: The curve $y=\dfrac{3}{x}$ and $x+y=4$.
\begin{list}{}{\setlength{\leftmargin}{3em}\setlength{\labelwidth}{2.5em}\setlength{\itemsep}{0pt}\setlength{\topsep}{2pt}}
\item[(A)] $1-2\ln 3$
\item[(B)] $3$
\item[(C)] $1-\ln 27$
\item[(D)] $4-\ln 27$
\end{list}

\bigskip\noindent\textbf{T8.}\ Find the area of the region with the following boundaries: In the first quadrant, bounded below by the $x$-axis and above by the curves of $y=2\sin x$ and $y=2\cos x$.
\begin{list}{}{\setlength{\leftmargin}{3em}\setlength{\labelwidth}{2.5em}\setlength{\itemsep}{0pt}\setlength{\topsep}{2pt}}
\item[(A)] $2(2-\sqrt{2})$
\item[(B)] $2(2+\sqrt{2})$
\item[(C)] $4$
\item[(D)] $2\sqrt{2}$
\end{list}

\bigskip\noindent\textbf{T9.}\ Find the area of the region with the following boundaries: Bounded above by the curve $y=2\sin x$ and below by $y=2\cos x$ from $x=\dfrac{\pi}{4}$ to $x=\dfrac{5\pi}{4}$.
\begin{list}{}{\setlength{\leftmargin}{3em}\setlength{\labelwidth}{2.5em}\setlength{\itemsep}{0pt}\setlength{\topsep}{2pt}}
\item[(A)] $4\sqrt{2}$
\item[(B)] $\dfrac{4}{\sqrt{2}}$
\item[(C)] $4(\sqrt{2}-1)$
\item[(D)] $4(\sqrt{2}+1)$
\end{list}

\bigskip\noindent\textbf{T10.}\ Find the area of the region with the following boundaries: The curve $y=\cot x$, the line $x=\dfrac{\pi}{3}$, and the $x$-axis.
\begin{list}{}{\setlength{\leftmargin}{3em}\setlength{\labelwidth}{2.5em}\setlength{\itemsep}{0pt}\setlength{\topsep}{2pt}}
\item[(A)] $\ln 2$
\item[(B)] $\dfrac{1}{2}\ln 3$
\item[(C)] $\ln 2-\dfrac{1}{2}\ln 3$
\item[(D)] $\dfrac{1}{2}\ln 2$
\end{list}

\bigskip\noindent\textbf{T11.}\ Find the area of the region with the following boundaries: The curve $y=x^3-3x^2-4x$ and the $x$-axis.
\begin{list}{}{\setlength{\leftmargin}{3em}\setlength{\labelwidth}{2.5em}\setlength{\itemsep}{0pt}\setlength{\topsep}{2pt}}
\item[(A)] $\dfrac{121}{4}$
\item[(B)] $\dfrac{131}{4}$
\item[(C)] $\dfrac{141}{4}$
\item[(D)] $\dfrac{151}{4}$
\end{list}

\bigskip\noindent\textbf{T12.}\ The total area bounded by the cubic $x=y^3-4y$ and the line $x=2y$ is equal to
\begin{list}{}{\setlength{\leftmargin}{3em}\setlength{\labelwidth}{2.5em}\setlength{\itemsep}{0pt}\setlength{\topsep}{2pt}}
\item[(A)] $9$
\item[(B)] $18$
\item[(C)] $6$
\item[(D)] $36$
\end{list}

\bigskip\noindent\textbf{T13.}\ The area bounded by $y=e^x$, $y=3$, and the $y$-axis is equal to
\begin{list}{}{\setlength{\leftmargin}{3em}\setlength{\labelwidth}{2.5em}\setlength{\itemsep}{0pt}\setlength{\topsep}{2pt}}
\item[(A)] $3-e$
\item[(B)] $e^2+1$
\item[(C)] $3\ln 3-2$
\item[(D)] $3\ln 3-4$
\end{list}

\bigskip\noindent\textbf{T14.}\ The area enclosed by the curve $y^2=x(4-x)$ is given by
\begin{list}{}{\setlength{\leftmargin}{3em}\setlength{\labelwidth}{2.5em}\setlength{\itemsep}{0pt}\setlength{\topsep}{2pt}}
\item[(A)] $\displaystyle 2\int_0^4 x\sqrt{4-x}\,dx$
\item[(B)] $\displaystyle 2\int_0^4 \sqrt{4x-x^2}\,dx$
\item[(C)] $\displaystyle 4\int_0^4 \sqrt{4x-x^2}\,dx$
\item[(D)] $\pi$
\end{list}

\bigskip\noindent\textbf{T15.}\ A rectangle has vertices $(0,0)$, $(c,0)$, $(c,c^2)$, and $(0,c^2)$, with $c>0$. The portion above the curve $y=x^2$ is shaded. What is the ratio of the area of the rectangle to its shaded area?
\begin{list}{}{\setlength{\leftmargin}{3em}\setlength{\labelwidth}{2.5em}\setlength{\itemsep}{0pt}\setlength{\topsep}{2pt}}
\item[(A)] $3:4$
\item[(B)] $5:4$
\item[(C)] $3:2$
\item[(D)] $3:1$
\end{list}

\bigskip\noindent\textbf{T16.}\ Find the volume generated by rotating the indicated region about the stated line: $y=x^2$, $x=3$, and $y=0$; about the $x$-axis.
\begin{list}{}{\setlength{\leftmargin}{3em}\setlength{\labelwidth}{2.5em}\setlength{\itemsep}{0pt}\setlength{\topsep}{2pt}}
\item[(A)] $\dfrac{243\pi}{5}$
\item[(B)] $27\pi$
\item[(C)] $\dfrac{81\pi}{5}$
\item[(D)] $\dfrac{162\pi}{7}$
\end{list}

\bigskip\noindent\textbf{T17.}\ Find the volume generated by rotating the indicated region about the stated line: $y=x^2$, $x=3$, and $y=0$; about the $y$-axis.
\begin{list}{}{\setlength{\leftmargin}{3em}\setlength{\labelwidth}{2.5em}\setlength{\itemsep}{0pt}\setlength{\topsep}{2pt}}
\item[(A)] $\dfrac{81\pi}{4}$
\item[(B)] $\dfrac{81\pi}{2}$
\item[(C)] $27\pi$
\item[(D)] $\dfrac{243\pi}{5}$
\end{list}

\bigskip\noindent\textbf{T18.}\ Find the volume generated by rotating the indicated region about the stated line: The first-quadrant region bounded by $y=x^2$, the $y$-axis, and $y=9$; about the $y$-axis.
\begin{list}{}{\setlength{\leftmargin}{3em}\setlength{\labelwidth}{2.5em}\setlength{\itemsep}{0pt}\setlength{\topsep}{2pt}}
\item[(A)] $9\pi$
\item[(B)] $\dfrac{81\pi}{2}$
\item[(C)] $\dfrac{243\pi}{2}$
\item[(D)] $27\pi$
\end{list}

\bigskip\noindent\textbf{T19.}\ Find the volume generated by rotating the indicated region about the stated line: $y=x^2$ and $y=9$; about the $x$-axis.
\begin{list}{}{\setlength{\leftmargin}{3em}\setlength{\labelwidth}{2.5em}\setlength{\itemsep}{0pt}\setlength{\topsep}{2pt}}
\item[(A)] $\dfrac{972\pi}{5}$
\item[(B)] $\dfrac{1944\pi}{5}$
\item[(C)] $\dfrac{486\pi}{5}$
\item[(D)] $\dfrac{243\pi}{10}$
\end{list}

\bigskip\noindent\textbf{T20.}\ Find the volume generated by rotating the indicated region about the stated line: $y=x^2$ and $y=9$; about the line $y=9$.
\begin{list}{}{\setlength{\leftmargin}{3em}\setlength{\labelwidth}{2.5em}\setlength{\itemsep}{0pt}\setlength{\topsep}{2pt}}
\item[(A)] $\dfrac{648\pi}{5}$
\item[(B)] $\dfrac{2592\pi}{5}$
\item[(C)] $\dfrac{1296\pi}{5}$
\item[(D)] $\dfrac{324\pi}{5}$
\end{list}

\bigskip\noindent\textbf{T21.}\ Find the volume generated by rotating the indicated region about the stated line: An arch of $y=\sin 2x$ and the $x$-axis; about the $x$-axis.
\begin{list}{}{\setlength{\leftmargin}{3em}\setlength{\labelwidth}{2.5em}\setlength{\itemsep}{0pt}\setlength{\topsep}{2pt}}
\item[(A)] $\dfrac{\pi}{4}\left(\pi-\dfrac{1}{4}\right)$
\item[(B)] $\dfrac{\pi^2}{2}$
\item[(C)] $\dfrac{\pi^2}{4}$
\item[(D)] $\dfrac{\pi^2}{8}$
\end{list}

\bigskip\noindent\textbf{T22.}\ Find the volume generated by rotating the indicated region about the stated line: A trapezoid with vertices at $(3,0)$, $(3,3)$, $(6,0)$, and $(6,6)$; about the $x$-axis.
\begin{list}{}{\setlength{\leftmargin}{3em}\setlength{\labelwidth}{2.5em}\setlength{\itemsep}{0pt}\setlength{\topsep}{2pt}}
\item[(A)] $\dfrac{56\pi}{3}$
\item[(B)] $63\pi$
\item[(C)] $189\pi$
\item[(D)] $126\pi$
\end{list}

\bigskip\noindent\textbf{T23.}\ The base of a solid is a circle of radius $b$ and is centered at the origin. Each cross section perpendicular to the $x$-axis is a square. The solid has volume
\begin{list}{}{\setlength{\leftmargin}{3em}\setlength{\labelwidth}{2.5em}\setlength{\itemsep}{0pt}\setlength{\topsep}{2pt}}
\item[(A)] $\dfrac{8}{3}b^3$
\item[(B)] $2\pi b^3$
\item[(C)] $4\pi b^3$
\item[(D)] $\dfrac{16}{3}b^3$
\end{list}

\bigskip\noindent\textbf{T24.}\ The base of a solid is the region bounded by the parabola $x^2=12y$ and the line $y=6$, and each plane section perpendicular to the $y$-axis is an equilateral triangle. The volume of the solid is
\begin{list}{}{\setlength{\leftmargin}{3em}\setlength{\labelwidth}{2.5em}\setlength{\itemsep}{0pt}\setlength{\topsep}{2pt}}
\item[(A)] $72\sqrt{3}$
\item[(B)] $216\sqrt{3}$
\item[(C)] $108\sqrt{3}$
\item[(D)] $432\sqrt{3}$
\end{list}

\bigskip\noindent\textbf{T25.}\ The base of a solid is the region bounded by $y=e^{-2x}$, the $x$-axis, the $y$-axis, and the line $x=1$. Each cross section perpendicular to the $x$-axis is a square. The volume of the solid is
\begin{list}{}{\setlength{\leftmargin}{3em}\setlength{\labelwidth}{2.5em}\setlength{\itemsep}{0pt}\setlength{\topsep}{2pt}}
\item[(A)] $e^4-1$
\item[(B)] $1-\dfrac{1}{e^4}$
\item[(C)] $\dfrac{e^4-1}{4}$
\item[(D)] $\dfrac{1}{4}\left(1-\dfrac{1}{e^4}\right)$
\end{list}

\bigskip\noindent\textbf{T26.}\ The length of the curve $y^2=x^3$ cut off by the line $x=9$ is
\begin{list}{}{\setlength{\leftmargin}{3em}\setlength{\labelwidth}{2.5em}\setlength{\itemsep}{0pt}\setlength{\topsep}{2pt}}
\item[(A)] $\dfrac{85\sqrt{85}-8}{27}$
\item[(B)] $\dfrac{85\sqrt{85}-1}{27}$
\item[(C)] $\dfrac{2}{27}(85\sqrt{85}-8)$
\item[(D)] $\dfrac{85\sqrt{85}}{27}$
\end{list}

\bigskip\noindent\textbf{T27.}\ The length of the curve $y=\ln(\sin x)$ from $x=\dfrac{\pi}{6}$ to $x=\dfrac{\pi}{4}$ equals
\begin{list}{}{\setlength{\leftmargin}{3em}\setlength{\labelwidth}{2.5em}\setlength{\itemsep}{0pt}\setlength{\topsep}{2pt}}
\item[(A)] $\displaystyle\int_{\pi/6}^{\pi/4} \sqrt{1+\tan^2 x}\,dx$
\item[(B)] $\displaystyle\int_{\pi/6}^{\pi/4} \sqrt{1+(\ln\sin x)^2}\,dx$
\item[(C)] $\displaystyle\int_{\pi/6}^{\pi/4} \sqrt{1+\dfrac{1}{\sin^2 x}}\,dx$
\item[(D)] $\displaystyle\int_{\pi/6}^{\pi/4} \sqrt{1+\cot^2 x}\,dx$
\end{list}

\bigskip\noindent\textbf{T28.}\ $\displaystyle\int_0^\infty e^{-2x}\,dx =$
\begin{list}{}{\setlength{\leftmargin}{3em}\setlength{\labelwidth}{2.5em}\setlength{\itemsep}{0pt}\setlength{\topsep}{2pt}}
\item[(A)] $1$
\item[(B)] $\dfrac{1}{2}$
\item[(C)] $-1$
\item[(D)] divergent
\end{list}

\bigskip\noindent\textbf{T29.}\ $\displaystyle\int_0^{e^2} \dfrac{du}{u} =$
\begin{list}{}{\setlength{\leftmargin}{3em}\setlength{\labelwidth}{2.5em}\setlength{\itemsep}{0pt}\setlength{\topsep}{2pt}}
\item[(A)] $2$
\item[(B)] $\dfrac{1}{e^2}$
\item[(C)] $-\dfrac{1}{e}$
\item[(D)] divergent
\end{list}

\bigskip\noindent\textbf{T30.}\ $\displaystyle\int_1^9 \dfrac{dt}{\sqrt[3]{t-1}} =$
\begin{list}{}{\setlength{\leftmargin}{3em}\setlength{\labelwidth}{2.5em}\setlength{\itemsep}{0pt}\setlength{\topsep}{2pt}}
\item[(A)] $\dfrac{2}{3}$
\item[(B)] $\dfrac{3}{2}$
\item[(C)] $6$
\item[(D)] divergent
\end{list}

\bigskip\noindent\textbf{T31.}\ $\displaystyle\int_{-4}^{12} \dfrac{dx}{(x-4)^{2/3}} =$
\begin{list}{}{\setlength{\leftmargin}{3em}\setlength{\labelwidth}{2.5em}\setlength{\itemsep}{0pt}\setlength{\topsep}{2pt}}
\item[(A)] $6$
\item[(B)] $0$
\item[(C)] $12$
\item[(D)] divergent
\end{list}

\bigskip\noindent\textbf{T32.}\ $\displaystyle\int_2^5 \dfrac{dx}{(x-4)^2} =$
\begin{list}{}{\setlength{\leftmargin}{3em}\setlength{\labelwidth}{2.5em}\setlength{\itemsep}{0pt}\setlength{\topsep}{2pt}}
\item[(A)] $2$
\item[(B)] $-2$
\item[(C)] $0$
\item[(D)] divergent
\end{list}

\bigskip\noindent\textbf{T33.}\ $\displaystyle\int_0^{\pi/2} \dfrac{2\sin x}{\sqrt{1-\cos x}}\,dx =$
\begin{list}{}{\setlength{\leftmargin}{3em}\setlength{\labelwidth}{2.5em}\setlength{\itemsep}{0pt}\setlength{\topsep}{2pt}}
\item[(A)] $-4$
\item[(B)] $1$
\item[(C)] $4$
\item[(D)] divergent
\end{list}

\bigskip\noindent\textbf{T34.}\ Find the area in the first quadrant under the curve $y=e^{-2x}$.
\begin{list}{}{\setlength{\leftmargin}{3em}\setlength{\labelwidth}{2.5em}\setlength{\itemsep}{0pt}\setlength{\topsep}{2pt}}
\item[(A)] $1$
\item[(B)] $\dfrac{1}{2}$
\item[(C)] $2$
\item[(D)] divergent
\end{list}

\bigskip\noindent\textbf{T35.}\ Find the area in the first quadrant under the curve $y=xe^{-x^2/2}$.
\begin{list}{}{\setlength{\leftmargin}{3em}\setlength{\labelwidth}{2.5em}\setlength{\itemsep}{0pt}\setlength{\topsep}{2pt}}
\item[(A)] $2$
\item[(B)] $\dfrac{2}{e}$
\item[(C)] $1$
\item[(D)] divergent
\end{list}

\bigskip\noindent\textbf{T36.}\ Find the area in the first quadrant above $y=2$, between the $y$-axis and the curve $y=\dfrac{3}{x}$.
\begin{list}{}{\setlength{\leftmargin}{3em}\setlength{\labelwidth}{2.5em}\setlength{\itemsep}{0pt}\setlength{\topsep}{2pt}}
\item[(A)] $1$
\item[(B)] $2$
\item[(C)] $\dfrac{1}{2}$
\item[(D)] divergent
\end{list}

\bigskip\noindent\textbf{T37.}\ Find the area between the curve $y=\dfrac{9}{1+x^2}$ and the $x$-axis.
\begin{list}{}{\setlength{\leftmargin}{3em}\setlength{\labelwidth}{2.5em}\setlength{\itemsep}{0pt}\setlength{\topsep}{2pt}}
\item[(A)] $\dfrac{9\pi}{4}$
\item[(B)] $\dfrac{9\pi}{2}$
\item[(C)] $9\pi$
\item[(D)] divergent
\end{list}

\bigskip\noindent\textbf{T38.}\ Find the area above the $x$-axis, between the curve $y=\dfrac{6}{\sqrt{1-x^2}}$ and its asymptotes.
\begin{list}{}{\setlength{\leftmargin}{3em}\setlength{\labelwidth}{2.5em}\setlength{\itemsep}{0pt}\setlength{\topsep}{2pt}}
\item[(A)] $3\pi$
\item[(B)] $6\pi$
\item[(C)] $9\pi$
\item[(D)] divergent
\end{list}

\bigskip\noindent\textbf{T39.}\ Find the volume of the solid generated when the region bounded above by $y=\dfrac{1}{x}$, at the left by $x=2$, and below by $y=0$ is rotated about the $x$-axis.
\begin{list}{}{\setlength{\leftmargin}{3em}\setlength{\labelwidth}{2.5em}\setlength{\itemsep}{0pt}\setlength{\topsep}{2pt}}
\item[(A)] $\dfrac{\pi}{4}$
\item[(B)] $\dfrac{\pi}{2}$
\item[(C)] $\pi$
\item[(D)] divergent
\end{list}

\bigskip\noindent\textbf{T40.}\ Find the volume of the solid generated when the first-quadrant region under $y=e^{-2x}$ is rotated about the $x$-axis.
\begin{list}{}{\setlength{\leftmargin}{3em}\setlength{\labelwidth}{2.5em}\setlength{\itemsep}{0pt}\setlength{\topsep}{2pt}}
\item[(A)] $\dfrac{\pi}{8}$
\item[(B)] $\dfrac{\pi}{4}$
\item[(C)] $\dfrac{\pi}{2}$
\item[(D)] divergent
\end{list}

\bigskip\noindent\textbf{TB1.}\ The area bounded by the parabola $y=3-x^2$ and the line $y=x-3$ is given by
\begin{list}{}{\setlength{\leftmargin}{3em}\setlength{\labelwidth}{2.5em}\setlength{\itemsep}{0pt}\setlength{\topsep}{2pt}}
\item[(A)] $\displaystyle\int_{-3}^{2} (6-x-x^2)\,dx$
\item[(B)] $\displaystyle\int_{-2}^{1} (3+x+x^2)\,dx$
\item[(C)] $\displaystyle\int_{-2}^{2} (6-x-x^2)\,dx$
\item[(D)] $\displaystyle 2\int_0^{\sqrt{3}} (3-x^2)\,dx + \int_{-3}^{2} (3-x)\,dx$
\end{list}

\bigskip\noindent\textbf{TB2.}\ For a continuous function $f$ on $[1,9]$, the table is $\begin{array}{c|ccccc}x&1&3&4&7&9\\\hline f(x)&2.0&5.2&8.6&10.4&13.5\end{array}$. Use a trapezoidal sum with four subintervals to approximate the area under the curve from $x=1$ to $x=9$, to two decimal places.
\begin{list}{}{\setlength{\leftmargin}{3em}\setlength{\labelwidth}{2.5em}\setlength{\itemsep}{0pt}\setlength{\topsep}{2pt}}
\item[(A)] $55.20$
\item[(B)] $61.30$
\item[(C)] $66.50$
\item[(D)] $70.80$
\end{list}

\bigskip\noindent\textbf{TB3.}\ The area $A$ enclosed by the three-leaved rose $r=2\cos 3\theta$ equals, to three decimal places,
\begin{list}{}{\setlength{\leftmargin}{3em}\setlength{\labelwidth}{2.5em}\setlength{\itemsep}{0pt}\setlength{\topsep}{2pt}}
\item[(A)] $0.785$
\item[(B)] $1.571$
\item[(C)] $3.142$
\item[(D)] $6.283$
\end{list}

\bigskip\noindent\textbf{TB4.}\ The area bounded by the small loop of the limaçon $r=1-2\cos\theta$ is given by the definite integral
\begin{list}{}{\setlength{\leftmargin}{3em}\setlength{\labelwidth}{2.5em}\setlength{\itemsep}{0pt}\setlength{\topsep}{2pt}}
\item[(A)] $\displaystyle\frac12\int_{-\pi/3}^{\pi/3}(1-2\cos\theta)^2\,d\theta$
\item[(B)] $\displaystyle\int_{\pi/3}^{5\pi/3} \left[\dfrac{1}{2}(1-2\cos\theta)\right]^2 d\theta$
\item[(C)] $\displaystyle\int_{2\pi/3}^{4\pi/3} \left[\dfrac{1}{2}(1-2\cos\theta)\right]^2 d\theta$
\item[(D)] $\displaystyle\int_{-2\pi/3}^{2\pi/3} \left[\dfrac{1}{2}(1-2\cos\theta)\right]^2 d\theta$
\end{list}

\bigskip\noindent\textbf{TB5.}\ Find the volume generated by rotating the indicated region about the stated line: $y=x^2$ and $y=9$; about the line $y=-2$.
\begin{list}{}{\setlength{\leftmargin}{3em}\setlength{\labelwidth}{2.5em}\setlength{\itemsep}{0pt}\setlength{\topsep}{2pt}}
\item[(A)] $\displaystyle 4\pi\int_{-2}^{9} (y+2)\sqrt{y}\,dy$
\item[(B)] $\displaystyle 2\pi\int_0^3 (9-x^2)^2\,dx$
\item[(C)] $\displaystyle \pi\int_{-3}^{3} (121-x^4)\,dx$
\item[(D)] $\displaystyle 2\pi\int_0^3 (117-4x^2-x^4)\,dx$
\end{list}

\bigskip\noindent\textbf{TB6.}\ The region bounded by $y=4x-x^2$ and $y=0$ is rotated about the $x$-axis. Which integral sets up its volume using the disk method with integration in $x$?
\begin{list}{}{\setlength{\leftmargin}{3em}\setlength{\labelwidth}{2.5em}\setlength{\itemsep}{0pt}\setlength{\topsep}{2pt}}
\item[(A)] $\displaystyle \pi\int_0^4 (16x^2+x^4)\,dx$
\item[(B)] $\displaystyle \pi\int_0^4 (4x-x^2)^2\,dx$
\item[(C)] $\displaystyle \pi\int_0^2 (4x-x^2)\,dx$
\item[(D)] $\displaystyle 2\pi\int_0^4 y\sqrt{16-4y}\,dy$
\end{list}

\bigskip\noindent\textbf{TB7.}\ Find the volume generated by rotating the indicated region about the stated line: $y=4x-x^2$ and $y=x$; about the $x$-axis.
\begin{list}{}{\setlength{\leftmargin}{3em}\setlength{\labelwidth}{2.5em}\setlength{\itemsep}{0pt}\setlength{\topsep}{2pt}}
\item[(A)] $\displaystyle \pi\int_0^2 \left[(4x-x^2)^2-x^2\right] dx$
\item[(B)] $\displaystyle \pi\int_0^3 (16x^2-8x^3)\,dx$
\item[(C)] $\displaystyle \pi\int_0^3 \left[(4x-x^2)^2-x^2\right] dx$
\item[(D)] $\displaystyle \pi\int_0^4 (3x-x^2)^2\,dx$
\end{list}

\bigskip\noindent\textbf{TB8.}\ Find the volume generated by rotating the indicated region about the stated line: $y=\ln x$, $y=0$, $x=e^2$; about the line $x=e^2$.
\begin{list}{}{\setlength{\leftmargin}{3em}\setlength{\labelwidth}{2.5em}\setlength{\itemsep}{0pt}\setlength{\topsep}{2pt}}
\item[(A)] $\displaystyle \pi\int_1^{e^2} (e^2-x)\ln x\,dx$
\item[(B)] $\displaystyle \pi\int_0^2 (e^2-e^y)^2\,dy$
\item[(C)] $\displaystyle 2\pi\int_1^{e^2} (e^2-\ln x)\,dx$
\item[(D)] $\displaystyle \pi\int_0^2 (e^4-2e^{y+1}+e^{2y})\,dy$
\end{list}

\bigskip\noindent\textbf{TB9.}\ A sphere of radius $r$ is divided into two parts by a plane at distance $h$ ($0<h<r$) from the center. The volume of the smaller part equals
\begin{list}{}{\setlength{\leftmargin}{3em}\setlength{\labelwidth}{2.5em}\setlength{\itemsep}{0pt}\setlength{\topsep}{2pt}}
\item[(A)] $\dfrac{\pi}{3}(2r^3+h^3-3r^2h)$
\item[(B)] $\dfrac{\pi h}{3}(3r^2-h^2)$
\item[(C)] $\dfrac{4}{3}\pi r^3+\dfrac{h^3}{3}-r^2h$
\item[(D)] $\dfrac{\pi}{3}(2r^3+3r^2h-h^3)$
\end{list}

\bigskip\noindent\textbf{TB10.}\ If the curves of $f(x)$ and $g(x)$ intersect for $x=a$ and $x=b$ and if $f(x)>g(x)>0$ for all $x$ on $(a,b)$, then the volume obtained when the region bounded by the curves is rotated about the $x$-axis is equal to
\begin{list}{}{\setlength{\leftmargin}{3em}\setlength{\labelwidth}{2.5em}\setlength{\itemsep}{0pt}\setlength{\topsep}{2pt}}
\item[(A)] $\displaystyle \pi\int_a^b (f(x))^2\,dx - \int_a^b (g(x))^2\,dx$
\item[(B)] $\displaystyle \pi\int_a^b [f(x)-g(x)]^2\,dx$
\item[(C)] $\displaystyle 2\pi\int_a^b x[f(x)-g(x)]\,dx$
\item[(D)] $\displaystyle \pi\int_a^b \left[(f(x))^2-(g(x))^2\right] dx$
\end{list}

\bigskip\noindent\textbf{TB11.}\ The length of one arch of the cycloid $x=2(t-\sin t)$, $y=2(1-\cos t)$ equals
\begin{list}{}{\setlength{\leftmargin}{3em}\setlength{\labelwidth}{2.5em}\setlength{\itemsep}{0pt}\setlength{\topsep}{2pt}}
\item[(A)] $\displaystyle\int_0^{\pi} 2\sqrt{1-\cos t}\,dt$
\item[(B)] $\displaystyle\int_0^{2\pi} 2\sqrt{\dfrac{1-\cos t}{2}}\,dt$
\item[(C)] $\displaystyle\int_0^{\pi} 2\sqrt{2-2\cos t}\,dt$
\item[(D)] $\displaystyle\int_0^{2\pi} 2\sqrt{2-2\cos t}\,dt$
\end{list}

\bigskip\noindent\textbf{TB12.}\ The length of the curve of the parabola $6x=y^2$ cut off by the line $x=3$ is given by the integral
\begin{list}{}{\setlength{\leftmargin}{3em}\setlength{\labelwidth}{2.5em}\setlength{\itemsep}{0pt}\setlength{\topsep}{2pt}}
\item[(A)] $\displaystyle\int_{-1}^{1} \sqrt{x^2+1}\,dx$
\item[(B)] $\displaystyle\dfrac{1}{3}\int_0^3 \sqrt{9+y^2}\,dy$
\item[(C)] $\displaystyle\int_{-1}^{1} \sqrt{1+x}\,dx$
\item[(D)] $\displaystyle\dfrac{2}{3}\int_0^{3\sqrt{2}} \sqrt{9+y^2}\,dy$
\end{list}

\bigskip\noindent\textbf{TB13.}\ The length of $x=e^t\cos t$, $y=e^t\sin t$ from $t=1$ to $t=2$ is equal to
\begin{list}{}{\setlength{\leftmargin}{3em}\setlength{\labelwidth}{2.5em}\setlength{\itemsep}{0pt}\setlength{\topsep}{2pt}}
\item[(A)] $6.605$
\item[(B)] $9.502$
\item[(C)] $11.318$
\item[(D)] $17.956$
\end{list}

\bigskip\noindent\textbf{TB14.}\ Which one of the following is an improper integral?
\begin{list}{}{\setlength{\leftmargin}{3em}\setlength{\labelwidth}{2.5em}\setlength{\itemsep}{0pt}\setlength{\topsep}{2pt}}
\item[(A)] $\displaystyle\int_0^3 \dfrac{dx}{\sqrt{x+2}}$
\item[(B)] $\displaystyle\int_{-1}^2 \dfrac{dx}{1+x^2}$
\item[(C)] $\displaystyle\int_0^3 \dfrac{x\,dx}{4-x^2}$
\item[(D)] $\displaystyle\int_0^{\pi/4} \dfrac{\sin x\,dx}{\cos^2 x}$
\end{list}

\bigskip\noindent\textbf{TB15.}\ Which one of the following improper integrals diverges?
\begin{list}{}{\setlength{\leftmargin}{3em}\setlength{\labelwidth}{2.5em}\setlength{\itemsep}{0pt}\setlength{\topsep}{2pt}}
\item[(A)] $\displaystyle\int_1^\infty \dfrac{dx}{x^3}$
\item[(B)] $\displaystyle\int_0^\infty \dfrac{dx}{e^{2x}}$
\item[(C)] $\displaystyle\int_{-1}^1\frac{dx}{\sqrt[4]{|x|}}$
\item[(D)] $\displaystyle\int_{-1}^1 \dfrac{dx}{x^3}$
\end{list}

\bigskip\noindent\textbf{TB16.}\ Which one of the following improper integrals diverges?
\begin{list}{}{\setlength{\leftmargin}{3em}\setlength{\labelwidth}{2.5em}\setlength{\itemsep}{0pt}\setlength{\topsep}{2pt}}
\item[(A)] $\displaystyle\int_0^\infty \dfrac{dx}{1+x^2}$
\item[(B)] $\displaystyle\int_0^1 \dfrac{dx}{x^{1/4}}$
\item[(C)] $\displaystyle\int_0^\infty \dfrac{dx}{e^x+3}$
\item[(D)] $\displaystyle\int_0^\infty \dfrac{dx}{x^{1/4}}$
\end{list}

\end{document}
